\section{Discussion}\label{sec:discussion}

The results separate two questions that current memory architectures answer with
one mechanism. A retriever answers \emph{what is relevant to the present
context?} Under a verification budget the agent must also answer \emph{which
inherited belief is most costly to leave unverified?} These need not induce the same ranking. \S\ref{sec:directs} shows the agent's scarce checking follows the
first question --- attention goes where the plan already points --- and
\S\ref{sec:content} shows the second also depends on the memory's own content:
a belief that states its constraint is checked substantially less often than a
length-matched belief that does not. We
note what this does \emph{not} show: no arm in any experiment states a
constraint that is false or stale, so whether a confidently-stated but outdated
constraint escapes checking --- the case the safety motivation turns on ---
remains untested.

Nothing here shows that the allocation we measure is \emph{wrong}. Every
constraint in every arm of every experiment is true, and declining to spend a
scarce credit re-deriving a true, already-stated constraint is locally sensible.
The concern is narrower and conditional: \emph{if} present relevance and future
expected loss diverge, a policy dominated by the former could under-check
beliefs whose later staleness would be costly. We do not establish that this
occurs. Testing it requires a constraint that is false or out of date, which we
do not run.

With that bound stated, verification scheduling is worth separating from
retrieval and from provenance storage as a control layer with its own inputs: relevance to the current plan,
loss if the belief is wrong, how many consolidation steps the record has passed
through, how far its scope has broadened, and when it was last checked. We do
not evaluate such a scheduler, and we could not identify a normatively correct
verification order from this data: that would require a corruption prior, a
distribution over future contexts and a loss model, none of which these
experiments supply. The contribution is the measurement that motivates building
one.

\section{Limitations}\label{sec:limits}

\textbf{Stochastic instability at the episode level.} Re-running
\StabilityN{} byte-identical prompts flips \StabilityFlip{} outcomes
(\StabilityFlipPct\%), ranging from \StabilityMin\% to \StabilityMax\% by model,
with \StabilityWorstModel{} least stable. Small effects here require replication before they are interpreted; that is why \S\ref{sec:content} reports primary and replication runs, and why we do not interpret contrasts of a few points.

\textbf{Model heterogeneity.} Magnitudes vary by an order of magnitude across
models (\S\ref{sec:content}). We avoid claims about language models in general.

\textbf{Geometry and scenarios.} Six inherited memories and a budget of one to
three is small against a production store, and the scenarios are synthetic. The allocation signature was reproduced across \SixModels{} models in two
worlds beyond the one the mechanism experiments use, but those mechanism
experiments use a single world.

\textbf{Prompting and structured output.} All results are obtained through a
fixed system prompt and a strict JSON schema. \S\ref{sec:directs} shows the main
causal result strengthens when the action field is removed, but the protocol is
not the only way to elicit these behaviours.

\textbf{The plan manipulation bundles three things.} \S\ref{sec:directs}
cannot separate working-plan state from the topical salience of naming the plan
or from directive framing. The randomised intervention effect stands; its
internal composition is not identified, and the discriminating cell was not
run.

\textbf{What is not identified.} The upstream and downstream halves of the chain
are identified separately and not end to end. Intent is measured after treatment
and is a candidate mediator; we compute no contrast on it and claim no
mediation. \S\ref{sec:content} does not identify a causal main effect of
register.

\textbf{Adaptive history.} Several interpretations reported in earlier versions
of this work were falsified by our own pre-registered tests and have been
withdrawn: the hedging effect of \S\ref{sec:content}, and the claim that
naturally consolidated records are harder to catch than the record they came
from. Pre-registrations, amendments and every episode file are in the
supplementary record.
